\chapter{Pharmacogenetic Testing}

\section*{What Is This Test?}
Pharmacogenetic testing analyzes germline genetic variants in genes that encode drug-metabolizing enzymes, drug transporters, and drug targets to predict an individual's response to specific medications. The purpose is personalized medicine: identifying patients at risk for drug toxicity (due to reduced metabolism and drug accumulation), therapeutic failure (due to ultrarapid metabolism and subtherapeutic drug levels), or adverse drug reactions prior to prescribing \cite{relling2015pharmacogenomics}. The most clinically validated pharmacogenetic tests include: CYP2C19 (clopidogrel activation, proton pump inhibitor metabolism, selective serotonin reuptake inhibitor metabolism), CYP2D6 (codeine/tramadol activation, tamoxifen activation, tricyclic antidepressant/nortriptyline metabolism, antipsychotic dosing), CYP2C9 (warfarin dosing, phenytoin metabolism), VKORC1 (warfarin sensitivity), TPMT and NUDT15 (thiopurine toxicity --- azathioprine, 6-mercaptopurine, thioguanine toxicity with myelosuppression risk), UGT1A1 (irinotecan toxicity --- severe neutropenia and diarrhea), DPYD (fluoropyrimidine toxicity --- 5-fluorouracil, capecitabine), SLCO1B1 (statin-induced myopathy --- simvastatin), HLA-B*57:01 (abacavir hypersensitivity syndrome), and HLA-B*15:02 and HLA-A*31:01 (carbamazepine-induced Stevens-Johnson syndrome/TEN). Pharmacogenetic results are reported as a predicted phenotype (e.g., poor metabolizer, intermediate metabolizer, normal metabolizer, rapid metabolizer, ultrarapid metabolizer) based on the combination of inherited alleles. Testing is performed by PCR-based genotyping of specific variant alleles or by NGS panels that include pharmacogenes.

\section*{Alternative Names}
Pharmacogenomics; PGx Testing; Drug-Gene Testing; Pharmacogenetic Panel; CYP450 Genotyping; Drug Metabolism Genetic Testing; Personalized Medicine Testing

\section*{Principle}
Pharmacogenetic testing uses PCR-based genotyping with allele-specific probes (TaqMan) or microarray-based platforms (e.g., DMET Plus, VeriDose, PharmacoScan) to detect specific single nucleotide polymorphisms (SNPs) and copy number variations (gene deletions and duplications) in pharmacogenes. For CYP2D6, the highly polymorphic gene locus includes gene deletions (*5), duplications/multiplications (*1xN, *2xN), and point mutations (*3, *4, *6, *10, *17, *41) that produce a range of enzyme activities from none (poor metabolizer, e.g., *4/*4) to markedly increased (ultrarapid metabolizer, e.g., *1/*1xN or *2/*2xN). For HLA alleles, high-resolution sequence-based typing is required (HLA-B*57:01 for abacavir, HLA-B*15:02 for carbamazepine in Asian populations). The Clinical Pharmacogenetics Implementation Consortium (CPIC) publishes standardized guidelines for translating genotype into phenotype and phenotype into specific drug prescribing recommendations \cite{relling2020cpic}. CPIC assigns level A (strongest evidence for clinical action) to gene-drug pairs including: TPMT/NUDT15 and thiopurines, CYP2C19 and clopidogrel, CYP2C9/VKORC1 and warfarin, CYP2D6 and codeine, DPYD and fluoropyrimidines.

\section*{History}
The term ``pharmacogenetics'' was coined by Friedrich Vogel in 1959. The discovery of the genetic basis of primaquine-induced hemolysis (G6PD deficiency) in the 1950s and succinylcholine apnea (pseudocholinesterase deficiency) were early milestones. In the 1980s--1990s, the CYP2D6 polymorphism was identified and characterized as the cause of the debrisoquine/sparteine oxidation polymorphism (poor metabolizer phenotype). The 2000s saw the emergence of CPIC (established in 2009), the Dutch Pharmacogenetics Working Group (DPWG), and the FDA pharmacogenomic labeling of drugs. In 2008, the FDA added a black box warning to abacavir recommending HLA-B*57:01 testing before prescription; this remains the gold standard example of pharmacogenetic testing preventing a life-threatening adverse reaction (hypersensitivity syndrome) \cite{mallal2008hlab5701}. In 2010, the FDA required a boxed warning on clopidogrel about reduced effectiveness in CYP2C19 poor metabolizers. In 2018, the FDA and CPIC recommended DPYD testing before fluoropyrimidine therapy.

\section*{Why This Test Is Ordered}
Pre-treatment testing to guide drug selection (clopidogrel vs. prasugrel/ticagrelor based on CYP2C19). Pre-treatment testing to guide initial drug dosing (warfarin based on CYP2C9 and VKORC1). Testing to predict risk of severe toxicity before initiating high-risk drugs (TPMT/NUDT15 before azathioprine/6-MP; DPYD before 5-FU/capecitabine; UGT1A1 before irinotecan; HLA-B*57:01 before abacavir) \cite{deenen2016upfront}. Testing to guide analgesic selection (CYP2D6 before codeine or tramadol). Testing to explain an unexpected adverse drug reaction or therapeutic failure potentially related to abnormal drug metabolism.

\section*{Primary vs Secondary Test}
This may be either a primary screening test performed preemptively before drug initiation (e.g., HLA-B*57:01 before abacavir is standard of care) or a secondary test to investigate an adverse drug reaction.

\section*{How This Test Is Performed}
Blood sample (EDTA tube) or buccal swab for DNA extraction. Genotyping by PCR or microarray. Results typically available in 3--10 days (or 1--3 days for stat/rapid pharmacogenetic testing).

\section*{What Preparation Is Needed}
No fasting or special preparation is required. Because the test analyzes constitutive (germline) DNA, results are not affected by drug dose, timing, or current medication use, so no drugs need to be held and no trough or peak timing applies. Blood is collected in an EDTA (purple-top) tube or, alternatively, a buccal swab is obtained. The ordering clinician should document the medications of interest and their indications so that the results can be translated into concrete prescribing recommendations.

\section*{Contraindications and Precautions}
There are no contraindications to blood sampling or buccal swab collection. Important interpretive cautions: (1) pharmacogenetic results are permanent and describe the patient's inherited metabolizer phenotype; they do not replace therapeutic drug monitoring and are not a substitute for clinical judgment in the setting of renal or hepatic impairment. (2) Results must be interpreted by a clinician or pharmacist familiar with CPIC guidelines, because genotype-to-phenotype translation is allele-specific and certain alleles (e.g., CYP2D6 duplications) require copy-number testing. (3) Ethnicity informs which alleles are relevant, and testing that screens only common alleles may miss rare or novel variants. (4) HLA testing must be high-resolution (sequence-based) for alleles such as HLA-B*57:01, and a negative result has a high but not absolute negative predictive value for abacavir hypersensitivity. (5) Genotype results carry implications for family members, and cascade testing and counseling should be offered when a clinically significant variant is found.

\section*{Risks}
Minimal risk. Buccal swab collection carries essentially no risk. Venipuncture carries slight pain or bruising at the puncture site, and rarely hematoma, infection, or vasovagal syncope.

\section*{What Happens After the Test}
Results are typically available in 3--10 days (1--3 days for rapid testing) and are reported as a predicted phenotype (e.g., poor, intermediate, normal, rapid, or ultrarapid metabolizer) with CPIC-based prescribing recommendations. A pharmacist or clinician reviews the results, adjusts the drug or dose (for example, selecting prasugrel or ticagrelor for a CYP2C19 poor metabolizer, or reducing thiopurine doses in TPMT deficiency), and discusses the implications with the patient. Confirmatory therapeutic drug monitoring is used when dose adjustment is based on phenotype, and relevant family members may be offered testing.

\section*{Understanding Results}

\subsection*{Poor Metabolizer (PM)}
Two nonfunctional alleles (e.g., CYP2C19 *2/*2). The enzyme activity is absent or profoundly reduced. Drug is NOT activated (prodrugs: clopidogrel, codeine, tamoxifen): therapy will be ineffective --- use an alternative agent. Drug is inactivated (active drugs: SSRIs, TCAs, proton pump inhibitors): risk of toxicity from drug accumulation --- reduce dose by 50--75\% or avoid. Examples: CYP2C19 PM + clopidogrel = 2--3-fold increased risk of stent thrombosis, MI, and death; switch to prasugrel or ticagrelor \cite{scott2013cyp2c19}. TPMT homozygous-deficient + azathioprine = profound myelosuppression (life-threatening neutropenia) at standard doses; reduce dose by 90\% or use alternative therapy \cite{relling2019thiopurine}. DPYD homozygous-deficient + 5-FU/capecitabine = catastrophic toxicity; avoid fluoropyrimidines entirely.

\subsection*{Intermediate Metabolizer (IM)}
One nonfunctional allele and one reduced-function allele, or two reduced-function alleles (e.g., CYP2C19 *1/*2). Enzyme activity is reduced. Adjust dose or consider alternative therapy. CYP2C19 IM + clopidogrel: 1.5--2-fold increased risk of adverse cardiovascular events; prasugrel or ticagrelor preferred.

\subsection*{Normal Metabolizer (NM)/Extensive Metabolizer (EM)}
Two fully functional alleles (e.g., CYP2D6 *1/*1). Standard drug dosing. Normal expected efficacy and toxicity.

\subsection*{Rapid (RM) or Ultrarapid Metabolizer (UM)}
Increased enzyme activity (gene duplications/multiplications, e.g., CYP2D6 *1/*2xN). Prodrugs: excessive conversion to active metabolite --- risk of toxicity (CYP2D6 UM + codeine: excessive conversion to morphine resulting in life-threatening respiratory depression/narcosis, especially in children and breastfeeding mothers) \cite{crews2014cyp2d6}. Active drugs: rapid elimination --- risk of therapeutic failure at standard doses (CYP2C19 UM + voriconazole: subtherapeutic antifungal levels). CYP2C19 UM + clopidogrel: increased conversion to active metabolite --- potentially enhanced platelet inhibition and bleeding risk, though clinical significance uncertain.

\subsection*{HLA Risk Allele Positive}
HLA-B*57:01 positive: Contraindication to abacavir --- abacavir should never be prescribed due to risk of multi-organ hypersensitivity syndrome (fever, rash, GI symptoms, respiratory failure, eosinophilia, death in <1--5\%). Negative predictive value of HLA-B*57:01 for abacavir hypersensitivity approaches 100\%. HLA-B*15:02 positive: Contraindication to carbamazepine in Asian populations due to risk of Stevens-Johnson syndrome/toxic epidermal necrolysis. HLA-A*31:01 positive: Increased risk of carbamazepine-induced DRESS syndrome.

\section*{Normal Results}
Normal metabolizer (NM) or extensive metabolizer (EM) phenotype for all tested pharmacogenes. The patient is expected to have normal drug metabolism and normal response to standard dosing. HLA negative for risk alleles. Results are reported as phenotype per CPIC terminology and specific prescribing recommendations (e.g., ``Consider alternative antiplatelet therapy'').

\section*{Follow-up Tests}
Drug-specific therapeutic drug monitoring (serum drug levels) if phenotype-predicted dosing is uncertain. Preemptive pharmacogenetic panel testing (multi-gene PGx panel) to inform future prescribing. Cascade testing of family members if a clinically significant pharmacogenetic variant is identified.

\section*{References}
\bibliographystyle{unsrtnat}
\bibliography{references}

\clearpage
